\section{Non-Functional Requirements}
\label{sec:non-functional-requirements}

Tab.~\ref{tab:nfr} presents the non-functional requirements that govern the
overall quality attributes of the system.
They were derived from the architectural goals stated in the problem statement
and from the operational constraints imposed by a researcher-run eye-tracking
experiment.
Priorities follow the same MoSCoW scheme \cite{clegg1994moscow} as Tab.~\ref{tab:fr}; not all
non-functional requirements are equally critical to the core thesis contribution.

\begingroup
\small
\begin{longtable}{%
  p{0.07\linewidth}
  p{0.18\linewidth}
  p{0.33\linewidth}
  p{0.25\linewidth}
  p{0.05\linewidth}
}
\caption{Non-functional requirements with rationale and MoSCoW priority.}
\label{tab:nfr} \\
\toprule
\textbf{ID} & \textbf{Quality attribute} & \textbf{Requirement} & \textbf{Rationale} & \textbf{Pri.} \\
\midrule
\endfirsthead

\multicolumn{5}{c}{\tablename~\thetable{} — continued} \\
\toprule
\textbf{ID} & \textbf{Quality attribute} & \textbf{Requirement} & \textbf{Rationale} & \textbf{Pri.} \\
\midrule
\endhead

\midrule
\multicolumn{5}{r}{\footnotesize Continued on next page} \\
\endfoot

\bottomrule
\endlastfoot

NFR1 & Modularity &
  The system shall maintain strictly separate bounded layers for sensing,
  eye-movement analysis, decision strategies, intervention execution,
  researcher interface, and session management.
  Each boundary shall be expressed as an interface in the application layer;
  concrete implementations shall be substitutable without changes to the core. &
  Motivated by the future research teams stakeholder
  (Sec.~\ref{sec:stakeholders}); the modularity claim must be demonstrable
  by substituting components at test time. & M \\
\midrule

NFR2 & Low Latency &
  The system shall keep the gaze-to-intervention latency, measured from a
  gaze sample to the dispatch of the resulting intervention, below
  100~ms, the threshold at which an interface response ceases to be
  perceived as instantaneous \cite{miller1968response, nielsen1993responsetime},
  and well within the duration of a single reading fixation
  (typically 200--300~ms) \cite{rayner1998eyemovements}.
  High-frequency gaze data shall be handled through a dedicated real-time
  channel, separate from general request-response communication. &
  Micro-interventions must be applied while the trigger condition is still
  present; a change delivered after the fixation that triggered it has
  ended would no longer match what the reader is looking at, invalidating
  the intervention rationale. The 100~ms ceiling is the conservative bound
  from interface-response perception, while the fixation timescale is the
  stricter reading-specific target. & M \\
\midrule

NFR3 & Usability &
  The system shall reduce experimenter cognitive load through a guided
  multi-step setup workflow, explicit readiness gating, and a unified
  control panel for all live session actions. &
  See Researcher actor (Sec.~\ref{sec:stakeholder-researcher});
  minimising operator cognitive load is identified there as a core
  requirement, not a peripheral usability enhancement. & S \\
\midrule

NFR4 & Reliability &
  The system shall fail gracefully on eye-tracker disconnection, invalid
  license, gaze-stream interruption, and external-provider disconnection:
  logging the failure, surfacing a status indicator to the researcher,
  and continuing in a degraded mode where possible. &
  An experiment session cannot be re-run if a crash causes data loss or
  disturbs the participant mid-reading. & M \\
\midrule

NFR5 & Reproducibility &
  All session parameters, calibration metrics, intervention settings,
  decision-strategy configuration, and system version information shall
  be captured in the session export so that a session can be understood
  and reproduced from that record alone. &
  See Data Consumer actor (Sec.~\ref{sec:stakeholder-consumer});
  their need for self-contained, complete session records is the direct
  source of this constraint. & M \\
\midrule

NFR6 & Extensibility &
  Adding new intervention modules, decision strategies, or eye-movement
  analysis algorithms shall require only additive code changes and shall
  not require modification of existing modules or the core layer.
  External processes shall be connectable as module providers without
  any changes to the codebase. &
  See External Module Provider actor (Sec.~\ref{sec:stakeholder-provider});
  the stable integration protocol requirement is motivated there. & M \\
\midrule

NFR7 & Recoverability &
  The system shall periodically persist session state to disk during an
  active experiment.
  An interrupted session shall be detectable and restorable on the next
  startup without manual intervention by the researcher. &
  Hardware and software faults during a live session with a human
  participant are disruptive; even partial data recovery preserves the
  experimental record. & M \\
\midrule

NFR8 & Hardware Independence &
  The system shall be fully operable without a physical eye tracker by
  using mouse-cursor simulation or another software sensing mode, with
  no code changes required. &
  Development, CI verification, and demonstrations must be possible on
  machines without Tobii hardware attached; this property also enables
  testing of all architectural layers in isolation. It is the
  quality-attribute counterpart of FR17.3, which mandates the same
  capability at the functional level. & M \\

\end{longtable}
\endgroup
